\documentclass[11pt]{article}
\usepackage[T1]{fontenc}
\usepackage{lmodern,amsmath,amssymb,amsthm,booktabs}
\usepackage[margin=1in]{geometry}
\usepackage[colorlinks=true,urlcolor=blue,citecolor=blue]{hyperref}
\newtheorem{proposition}{Proposition}
\newcommand{\rem}{\operatorname{rem}}
\newcommand{\eq}{\operatorname{eq}}
\newcommand{\lean}[1]{\nolinkurl{#1}}
\makeatletter\g@addto@macro\UrlBreaks{\do\_}\makeatother
\emergencystretch=2em
\title{Growing prefixes in a fixed remainder program}
\author{Research note accompanying the uniform-process construction}
\date{October 16, 2026}
\begin{document}
\maketitle

\begin{abstract}
The fixed-program family does not yet give one sublogarithmic remainder
program. We investigate recomputation, destructive traversal, and recursive
reuse of a growing prefix. A concrete finite program of maximum arity three
can select arbitrarily deep prefix operands, but the tested computations pay
for repeated regeneration. We formalize a limitation of the corresponding
cost accounting, and describe a structural direction based on the interfaces
between recursive calls. No sublogarithmic algorithm or algorithmic lower
bound is established here.
\end{abstract}

\section{The target and the operational model}

The desired statement is the existence of one finite nondeterministic
recursive program over $(\mathbb N,\rem,\eq_0)$ computing the gcd with
\[
\max_{1\leq b\leq a} c_P(a,b)=o(\log a),
\]
where $c_P$ is the minimum remainder count of a convergent computation.
Every convergent computation must be correct, and at least one must exist
on every input. These are the conventions of \cite[\S2E]{book}.
The existing theorem chooses a different fixed program for each desired
coefficient; its generator arity is $4^p+5$ and its coefficient is at most
$3/p$. The quantifiers cannot simply be exchanged.

The module \lean{ShortGCD.RegisterSmallStep} now formalizes the specialized
two-stack recursive machine. It proves exact-cost equivalence with finite
evaluation trees, including right-to-left argument evaluation, separate zero
tests and branches, finite nondeterministic body choice, and a positive
divisor requirement. Transition locality is also checked. A saved caller
context survives a child computation, but the child cannot inspect arbitrary
saved values. The book itself distinguishes the richer structure used to
implement the stacks from the original structure \cite[p.~78]{book}.
Machine stacks therefore do not supply list primitives to the source program.

\section{A single program with unbounded prefix selection}

Write $a=r_0,b=r_1,r_2,\ldots$ for the Euclidean sequence. The fixed equations
\begin{align*}
\operatorname{prefix}(a,b)&=a,\\
\operatorname{prefix}(a,b)&=b,\\
\operatorname{prefix}(a,b)&=\operatorname{prefix}(b,\rem(a,b))
\end{align*}
can select $r_i$ with exactly $\max(0,i-1)$ remainder calls, whenever the
necessary divisors are positive. Thus arbitrary selection depth itself is
available in one finite program.

Define the four alternatives
\[
\operatorname{near}(u,v)=u\;\mid\;v\;\mid\;\rem(u,v)\;\mid\;\rem(v,u)
\]
and the single equation
\[
\operatorname{shrink}(a,b,m)=\operatorname{near}
 \bigl(\rem(\operatorname{prefix}(a,b),m),
       \rem(\operatorname{prefix}(a,b),m)\bigr).
\]
Finally let
\begin{align*}
\operatorname{generate}(a,b,m)&=m,\\
\operatorname{generate}(a,b,m)&=
 \operatorname{generate}(a,b,\operatorname{shrink}(a,b,m)).
\end{align*}
Use the earlier verification wrapper: if $b=0$, return $a$; otherwise
generate a candidate starting with $m=b$ and return it only if it is positive
and divides both original inputs. Rejection enters a self-loop. All tests
are implemented using $\eq_0$ and remainder.

There are seven internal functions including the entry function and rejection
loop, all of arity at most three. There are no numeric constants, list
operations, or runtime indices in the program. Its syntax is independent of
every input and of the planner's prefix length.

Every common divisor of $a,b$ divides every generated candidate, by induction
on evaluations. The wrapper therefore makes every convergent computation
correct. A successful computation always exists: if $b=\gcd(a,b)$, stop
immediately; otherwise choose the gcd as both prefix operands. Their residues
modulo $b$ equal the gcd, so a projection alternative of near returns it.
For $a\geq b>0$, if Euclid takes $E$ calls to make the gcd available, this
witness costs at most $2E+4$ calls, including verification. This proves only
an $O(\log a)$ upper bound.

\paragraph{Executable check.}
\lean{rebuilding_program.py} compiles these equations to the existing finite
AST. The independently defined right-to-left interpreter executes the same
compiled object for every test and planning depth. The planner chooses close
residues among the first $h+1$ prefix values; it uses unrestricted arithmetic
offline and provides only the nondeterministic choices to the interpreter.
Selecting prefix positions $i,j$ costs
\[
\max(0,i-1)+\max(0,j-1)+2+e,
\qquad e\in\{0,1\},
\]
where the two calls produce residues and $e$ records a remainder in near.
Both argument occurrences are evaluated, even when the selected values agree.

The suite checks 270 runs and 270 additional scaled replays with identical
choices. It includes all pairs below nine, unordered and zero inputs, larger
random and scaled pairs, and a Fibonacci pair. On one seeded 2,048-bit input:
\begin{center}
\begin{tabular}{rrrrr}
\toprule
Planning $h$ & Contractions & Regeneration calls & Local calls & Total calls\\
\midrule
2 & 538 & 373 & 1,566 & 1,941\\
5 & 349 & 1,056 & 1,042 & 2,100\\
45 & 169 & 7,359 & 506 & 7,867\\
\bottomrule
\end{tabular}
\end{center}
Totals include two final verification calls. Increasing the available prefix
reduces contractions on this example but increases operand regeneration.
These are costs of particular selected runs, not minimum costs. They do not
prove that this program has no sublogarithmic successful computations or
that a better planner cannot find them.

\section{What the prefix estimate can certify without reuse}

For a target contraction factor $2^p$, the existing prefix lemma uses
$h=4^p+1$ ordered pairs. Building that prefix from its initial pair takes
$h-1=4^p$ Euclidean calls. The lemma guarantees at least $p$ bits of decrease.
Suppose an analysis budgets a fresh prefix for each contraction, discards it
afterwards, and accounts for progress only by these guaranteed factors.
It may choose a different positive integer $p_i$ at every stage.

\begin{proposition}[Restricted accounting bound]
For any such finite list $p_1,\ldots,p_J$,
\[
\sum_{i=1}^J4^{p_i}\ \geq\ 4\sum_{i=1}^Jp_i.
\]
Consequently, if this accounting justifies a target decrease of $n$ bits
solely by $n\leq\sum_i p_i$, its charged prefix budget is at least $4n$.
\end{proposition}
\begin{proof}
For integers $p\geq1$, $4p\leq4^p$, by induction. Sum these inequalities
and use the assumed target inequality.
\end{proof}

\lean{PrefixPolicies.lean} checks both inequalities. This is a lower bound
on the budget of this proof policy, not on actual executions. In particular,
the hypothesis $n\leq\sum_i p_i$ is not necessary for an actual successful
run: its contractions may be much stronger than their guaranteed factors.
The prototype also reconstructs only its selected operands, so it need not
pay the full fresh-prefix charge at every stage. The proposition cannot be
applied to its measured costs as an algorithmic lower bound.

Two related attempts have the same limitation at the level of this estimate.
Recomputing all prefix residues at every modulus replaces three calls per
contraction by a budget growing with the prefix. Consuming disjoint prefix
blocks while unwinding avoids repeated traversal, but spends a new block for
each application of the residue-count lemma. For blocks with $h_i=4^{p_i}+1$
ordered pairs, a consecutive Euclidean prefix containing all these blocks
takes $\sum_i h_i-1\geq\sum_i4^{p_i}$ calls when $J\geq1$.
Neither estimate yields a sublinear budget for a fixed positive fraction of
the input's bit length. This leaves open reuse across overlapping blocks,
stronger residue information, and schedules that exploit larger contractions.

\section{A structural direction: bounded interfaces between calls}

The following observation concerns computation occurrences, not numerical
values identified whenever they happen to agree. It suggests a way to state
the remaining issue more precisely. It is a mathematical argument, not a
Lean theorem or a result claimed to be new.

For a terminating evaluation, form its data-dependence graph. Give every
primitive result and every local expression result an occurrence vertex.
Give each call fresh formal-input ports and a return port, linked by copying
edges to its actual arguments and its returned expression. Repeated uses of
one formal input share that port. Keep separate vertices for unrelated
occurrences even if their numerical values are equal. Consider the underlying
undirected graph, optionally adding edges that make each primitive's operands
and result a clique.

\begin{proposition}[Call-tree decomposition]
For each fixed finite first-order recursive program $P$ without shared mutable
state, these occurrence graphs have treewidth bounded by a constant depending
only on $P$, independently of the input and recursion depth.
\end{proposition}
\begin{proof}
Use the tree of dynamic function invocations, adding the head evaluation as
its root. For each invocation, put in its bag its formal-input ports, the
results of all subterm occurrences executed in its own body, its return port,
and the formal-input and return ports of its immediate child invocations.
The number of such occurrences and child calls is bounded by the finite body
syntax, and each child has fixed arity. Taking the maximum over the finitely
many bodies and the head gives a bound $C(P)$ on bag size.

Every primitive edge lies in the bag of the invocation executing it. Edges
for argument and return copying lie in the caller or callee bag, and each
interface port lies in both adjacent bags. Other local vertices lie only in
their own bag. Thus the bags containing any vertex form a connected subtree,
and every edge is covered. This is a tree decomposition of width at most
$C(P)-1$.
\end{proof}

Copy vertices may be removed by contracting copying edges, and unused tests
and vertices may be deleted. These graph-minor operations preserve the
existence of a bounded-width decomposition. Thus a small-cost execution gives
a remainder derivation with a bounded-width occurrence graph as well.
This argument does not allow the unrestricted identification of equal
numerical values: that quotient need not preserve treewidth. Existing
certificate files store numerical values, so measuring the width of their
value-identified graphs would not by itself test this observation.

This leads to two concrete research questions.
\begin{enumerate}
\item Can the short prefix construction be organized into bounded-width
occurrence graphs, with a decomposition and execution schedule realizable by
one fixed program? The residue lemma supplies suitable operand pairs for each
modulus but gives no compatibility between their choices at different moduli.
Bounded treewidth alone is not asserted sufficient for such a compilation.
\item For every fixed width $t$, is there a $c_t>0$ and an infinite family of
inputs on which every remainder derivation of the gcd with occurrence
treewidth at most $t$ needs at least $c_t\log a$ calls? Such a theorem would
imply the nondeterministic Main Conjecture using the decomposition above.
No such lower bound is proved here. It would use a restriction on program
computations beyond the unrestricted homomorphism method.
\end{enumerate}

Treewidth restrictions already play a role in circuit complexity; see, for
example, Jansen and Sarma \cite{js}. Their balancing results do not give a
remainder-program speedup or the proposed lower bound. The immediate next
test is to instrument the existing interpreter with occurrence identities,
construct its invocation bags, and independently check the decomposition.
That would make the proposed restriction concrete without assuming that
recursion depth bounds storage or that equal values are interchangeable
across inaccessible frames.

\section{Reproduction and status}
From the repository root:
\begin{verbatim}
lake build
lake env lean research/RegisterProgram.lean
lake env lean research/PrefixPolicies.lean
python3 research/rebuilding_program.py
python3 -O research/rebuilding_program.py
\end{verbatim}
The machine correspondence is part of the normal library build and axiom
audit. The restricted budget lemma has its own audit in
\lean{prefix-policy-axioms.txt}. Prototype results are recorded in
\lean{rebuilding-validation.json}. The prototype and the call-tree
decomposition are not formalized in Lean. The established main-paper bounds
are unchanged; one fixed sublogarithmic program remains open.

\begin{thebibliography}{2}
\bibitem{book}
Yiannis N. Moschovakis.
\emph{Abstract Recursion and Intrinsic Complexity}.
Lecture Notes in Logic 48, Cambridge University Press / ASL, 2019.
\url{https://www.math.ucla.edu/~ynm/lectures/aricplace/c-aric.pdf}.
\bibitem{js}
Maurice Jansen and Jayalal Sarma M.N.
\emph{Balancing Bounded Treewidth Circuits}, 2009.
\url{https://arxiv.org/abs/0910.1427}.
\end{thebibliography}
\end{document}
